\par\Needspace{6\baselineskip}\medskip
\begin{center}\hypertarget{native-192-article143}{}\pdfbookmark[2]{143.}{article-143}143.\end{center}
Sobald der Winkel $z$ gefunden ist, erhält man $r'$ unmittelbar aus der Gleichung
\[
r'=\frac{R'\sin\delta'}{\sin z}\mathpunct{\cdot}
\]
Ferner folgt aus den Gleichungen $P=\dfrac{n''}{n}$ und III
\[
\frac{n'r'}{n}=\frac{(P+a)R'\sin\delta'}{b\sin(z-\sigma)}
\]
\[
\frac{n'r'}{n''}=\frac{1}{P}\cdot\frac{n'r'}{n}\mathpunct{\cdot}
\]

\SourceLeaf{193}Wir behandeln nun die Formeln, nach denen die Lage der Punkte $C,C''$ aus der Lage des Punktes $C'$ zu bestimmen ist, so, dass ihre allgemeine Gültigkeit auch für die in Abb. 4 nicht dargestellten Fälle sofort ersichtlich wird. Dazu bemerken wir, dass der Sinus des Abstands des Punktes $C'$ vom Großkreis $CB$ (der Abstand wird auf der oberen Seite positiv, auf der unteren negativ genommen) gleich dem Produkt aus $\sin\epsilon''$ und dem Sinus des in Vorwärtsrichtung gemessenen Abstands des Punktes $C'$ von $D''$ ist, also $=-\sin\epsilon''\sin C'D''=-\sin\epsilon''\sin(z+A'D''-\delta')$. Ebenso ist der Sinus des Abstands des Punktes $C''$ von demselben Großkreis $=-\sin\epsilon'\sin C''D'$. Diese Sinuswerte verhalten sich offenbar wie $\sin CC'$ zu $\sin CC''$ oder wie $\dfrac{n''}{rr'}$ zu $\dfrac{n'}{rr''}\mathpunct{\raisebox{.45ex}{,}}$ oder wie $n''r''$ zu $n'r'$. Setzt man also $C''D'=\zeta''$, so erhält man
\[
\tag*{V.}r''\sin\zeta''=\frac{n'r'}{n''}\cdot\frac{\sin\epsilon''}{\sin\epsilon'}\sin(z+A'D''-\delta').
\]
Ganz entsprechend ergibt sich mit $CD'=\zeta$
\[
\tag*{VI.}r\sin\zeta=\frac{n'r'}{n}\cdot\frac{\sin\epsilon}{\sin\epsilon'}\sin(z+A'D-\delta')
\]
\[
\tag*{VII.}r\sin(\zeta+AD''-AD')=r''P\cdot\frac{\sin\epsilon}{\sin\epsilon''}\sin(\zeta''+A''D-A''D').
\]
Verbindet man die Gleichungen V und VI mit den folgenden, aus Artikel 139 übernommenen Gleichungen,
\[
\tag*{VIII.}r''\sin(\zeta''-A''D'+\delta'')=R''\sin\delta''
\]
\[
\tag*{IX.}r\sin(\zeta-AD'+\delta)=R\sin\delta,
\]
so lassen sich die Größen $\zeta,\zeta'',r,r''$ nach der Regel in Artikel 78 ableiten. Um die Rechnung zu erleichtern, geben wir die Formeln hier ausdrücklich an. Man setze
\[
\tag*{17)}\frac{R\sin\delta}{\sin(AD'-\delta)}=\varkappa
\]
\[
\tag*{18)}\frac{R''\sin\delta''}{\sin(A''D'-\delta'')}=\varkappa''
\]
\[
\tag*{19)}\frac{\cos(AD'-\delta)}{R\sin\delta}=\lambda
\]
\[
\tag*{20)}\frac{\cos(A''D'-\delta'')}{R''\sin\delta''}=\lambda''.
\]
Die Berechnung dieser Größen, genauer ihrer Logarithmen, ist weiterhin unabhängig von $P$ und $Q$ und bildet die \emph{fünfte} und letzte der gewissermaßen vorbereitenden Rechnungen. Sie lässt sich bequem gleich mit der Berechnung von $a,b$, also mit der vierten Aufgabe, erledigen, denn es gilt $a=\dfrac{\varkappa}{\varkappa''}\mathpunct{\cdot}$ — Setzt man \SourceLeaf{194}dann
\[
\begin{aligned}
\frac{n'r'}{n}\cdot\frac{\sin\epsilon}{\sin\epsilon'}\cdot\sin(z+A'D-\delta')&=p\\[.5ex]
\frac{n'r'}{n''}\cdot\frac{\sin\epsilon''}{\sin\epsilon'}\cdot\sin(z+A'D''-\delta')&=p''\\[.5ex]
\varkappa(\lambda p-1)&=q\\[.5ex]
\varkappa''(\lambda''p''-1)&=q''
\end{aligned}
\]
so erhält man $\zeta$ und $r$ aus $r\sin\zeta=p$, $r\cos\zeta=q$ sowie $\zeta''$ und $r''$ aus $r''\sin\zeta''=p''$, $r''\cos\zeta''=q''$. Bei der Bestimmung von $\zeta$ und $\zeta''$ kann hier keine Mehrdeutigkeit auftreten, weil $r$ und $r''$ notwendig positiv sein müssen. Die abgeschlossene Rechnung kann man nach Belieben mit Gleichung VII überprüfen.

In zwei Fällen muss man jedoch anders vorgehen. Fällt der Punkt $D'$ mit $B$ zusammen oder liegt er ihm auf der Kugel gegenüber, gilt also $AD'-\delta=0$ oder $=180^0$, so müssen die Gleichungen VI und IX identisch sein. Dann ergäbe sich $\varkappa=\infty$, $\lambda p-1=0$, und $q$ wäre folglich unbestimmt. In diesem Fall bestimmt man $\zeta''$ und $r''$ wie angegeben. Anschließend muss man $\zeta$ und $r$ aus der Verbindung von Gleichung VII mit VI oder IX ableiten. Die Formeln selbst, die man Artikel 78 entnehmen kann, schreiben wir hier nicht noch einmal hin. Wir bemerken nur, dass dasselbe Verfahren auch dann vorzuziehen ist, wenn $AD'-\delta$ weder $=0$ noch $=180^0$, wohl aber ein sehr kleiner Bogen ist, weil das vorige Verfahren dann die nötige Genauigkeit nicht zulässt. Je nachdem, ob $\sin(AD''-AD')$ größer oder kleiner ist als $\sin(AD''-\delta)$, verbindet man Gleichung VII mit VI oder mit IX.

Fällt ebenso der Punkt $D'$ oder sein Gegenpunkt mit $B''$ zusammen oder liegt er nur wenig davon entfernt, so wäre die Bestimmung von $\zeta'',r''$ nach dem vorigen Verfahren entweder unmöglich oder wenig zuverlässig. Dann bestimmt man $\zeta$ und $r$ nach diesem Verfahren und leitet anschließend $\zeta''$ und $r''$ aus der Verbindung von Gleichung VII mit V oder VIII ab, je nachdem, ob $\sin(A''D-A''D')$ größer oder kleiner ist als $\sin(A''D-\delta'')$. Es ist jedoch nicht zu befürchten, dass \emph{gleichzeitig} $D'$ mit den Punkten $B,B''$ oder ihren Gegenpunkten zusammenfällt oder nur wenig von ihnen entfernt liegt. Denn den Fall, dass $B$ mit $B''$ zusammenfällt oder sehr nahe bei ihm liegt, haben wir bereits oben in Artikel 138 von unserer Untersuchung ausgeschlossen.
